\documentclass[11pt, a4paper]{article}

\usepackage[T1]{fontenc}
\usepackage[utf8]{inputenc}
\usepackage{tgpagella}
\usepackage[spanish, es-noshorthands]{babel}
\usepackage{amsmath, amssymb, bm}
\usepackage{tabularx}
\usepackage{booktabs}
\usepackage[most]{tcolorbox}
\usepackage[margin=2.5cm]{geometry}
\usepackage{ragged2e}
\usepackage{fancyhdr}

\pagestyle{fancy}
\fancyhf{}
\setlength{\headheight}{16pt}
\lhead{\textbf{UCB Sede Tarija} | Ing. Mecatrónica}
\chead{}
\rhead{IMT-121 Circuitos Electrónicos I | Reconciliación Hito 2}
\lfoot{Prof: M.Sc. Ing. Bernardo Quiroga Turdera}
\rfoot{Página \thepage}
\renewcommand{\headrulewidth}{0.4pt}

\definecolor{ucbblue}{RGB}{0, 51, 102}

\begin{document}

\begin{tcolorbox}[colback=ucbblue!5!white, colframe=ucbblue, title=\textbf{Matriz de Reconciliación de Evidencias (EC2)}]
    \centering \textbf{Consolidación de Calificaciones y Rutas de Acción}
\end{tcolorbox}

\noindent \textit{Uso Exclusivo Docente. Ponderación oficial aprobada: Examen Individual (50\%), Componente Práctico (50\% - integra trabajo en clase y Reporte IEEE A+B).}

\vspace{0.3cm}
\renewcommand{\arraystretch}{2} % Aumentado para que los renglones sean cómodos al escribir a mano
\noindent
% Ajuste proporcional exacto de las 8 columnas para evitar desbordamientos.
% La suma de los \hsize (2.0 + 0.6 + 0.8 + 0.8 + 0.7 + 0.9 + 1.5 + 0.7) es exactamente 8.
\begin{tabularx}{\textwidth}{
  >{\RaggedRight\arraybackslash\hsize=2.0\hsize}X 
  >{\centering\arraybackslash\hsize=0.6\hsize}X 
  >{\centering\arraybackslash\hsize=0.8\hsize}X 
  >{\centering\arraybackslash\hsize=0.8\hsize}X 
  >{\centering\arraybackslash\hsize=0.7\hsize}X 
  >{\centering\arraybackslash\hsize=0.9\hsize}X 
  >{\RaggedRight\arraybackslash\hsize=1.5\hsize}X 
  >{\centering\arraybackslash\hsize=0.7\hsize}X}
\toprule
\textbf{Estudiante} & \textbf{Grupo} & \textbf{Examen (50\%)} & \textbf{Lab. (50\%)} & \textbf{Total Prov.} & \textbf{Estado Evid.} & \textbf{Acción Req.} & \textbf{Ruta} \\ 
\midrule
 & & & & & & & \\ \midrule
 & & & & & & & \\ \midrule
 & & & & & & & \\ \midrule
 & & & & & & & \\ \midrule
 & & & & & & & \\ \midrule
 & & & & & & & \\ \midrule
 & & & & & & & \\ 
\bottomrule
\end{tabularx}

\vspace{0.5cm}
\textbf{Claves de Estado de Evidencia:}
\begin{itemize}
    \item \textbf{CONSISTENTE:} Examen, reporte y desempeño en banco guardan relación lógica. (Ruta A).
    \item \textbf{DISCREPANTE:} Diferencia masiva entre el desempeño individual (examen) y el trabajo entregado por el equipo (reporte). Activa Microdefensa (Ruta B).
\end{itemize}

\textbf{Definición de Ruta de la Sesión:}
\begin{itemize}
    \item \textbf{PATH A (Mayoría Consistente):} Síntesis final EC2 y puente conceptual a transitorios (EC3).
    \item \textbf{PATH B (Inconsistencias Selectivas):} Asignación de tarea de corrección al grupo mientras se ejecutan microdefensas 1-a-1.
    \item \textbf{PATH C (Dificultad Generalizada):} La cohorte falló un mismo núcleo (ej. Thévenin). Se activa remediación inmediata acotada (Artefacto 5).
\end{itemize}

\end{document}
